Branchless programming is common folklore for data-dependent operations such as sign-dependent
variable shifts ($v \ll s$ if $s>0$, $v \gg |s|$ otherwise). This paper examines whether the
folklore holds on the Cortex-M0+, an in-order microcontroller core without branch predictors,
by comparing a branch baseline, a branchless C implementation, and a hand-written assembly
implementation on real RP2040 hardware. The branchless C version was 1.2$\times$ slower than
the plain branch, because the C standard forbids the compiler from exploiting the ARM barrel
shifter's saturation semantics (shift amounts $\geq$ 32 yield zero) and the generated code spends
14 instructions per element. A hand-written assembly loop that delegates sign selection to the
shifter saturation needs only 8 instructions per element and runs 1.5$\times$ faster than the
branch baseline, with measured cycles matching the analytical model within 2\%. We analyze the
gap between compiler-generated and hand-written code, attribute it to language-rule constraints,
destructive two-operand encoding, and register allocation, and discuss a direction in which
large language models predict such compiler-pass limitations and synthesize new assembly
kernels or compiler passes, as demonstrated by the LLM-driven feedback loop used throughout
this study.
